\begin{center}
\textit{Allen Charakteren, außer den in Art.\ $263,264$ als unmöglich erkannten, entsprechen tatsächlich Geschlechter.}
\end{center}

\Needspace{8\baselineskip}
\Article{287}
Zum Problem des vorigen Artikels fügen wir noch die folgenden Bemerkungen hinzu.

I.\ Geht die Form $F$ durch die Substitution $p,p',p'',p''';\ q,q',q'',q'''$ in das Produkt
der beiden Formen $(h,i,k)$, $(h',i',k')$ über (beide, wie wir stets
\SourceLeaf{342}
voraussetzen, im eigentlichen Sinn genommen), so erhält man die folgenden
Gleichungen, die sich leicht aus Schluss 3 in Art.\ 235 ableiten lassen:
\[
\begin{gathered}
p''hn'-p'h'n-p(in'-i'n)=0\\
(p''-p')(in'+i'n)-p(kn'-k'n)+p'''(hn'-h'n)=0\\
p'kn'-p''k'n-p'''(in'-i'n)=0
\end{gathered}
\]
sowie drei weitere, die daraus durch Vertauschen der Zahlen $p$,
$p'$, $p''$, $p'''$ mit $q$, $q'$, $q''$, $q'''$ entstehen.
Dabei sind $n$, $n'$ die positiven Quadratwurzeln der Quotienten,
die man erhält, wenn man die Determinanten der Formen $(h,i,k)$, $(h',i',k')$
durch die Determinante der Form $F$ teilt. Sind also diese Formen
identisch, oder gilt $n=n'$, $h=h'$, $i=i'$, $k=k'$, so gehen jene
Gleichungen in die folgenden über:
\[
(p''-p')hn=0,\quad (p''-p')in=0,\quad (p''-p')kn=0
\]
Daraus folgt \textit{notwendig} $p'=p''$, und auf genau dieselbe Weise
$q'=q''$. --- Gibt man also den Formen $(h,i,k)$, $(h',i',k')$ \textit{dieselben}
Unbestimmten $t$, $u$ und bezeichnet die Unbestimmten der Form
$F$ mit $T$, $U$, so geht $F$ durch die folgende Substitution über:
\[
T=ptt+2p'tu+p'''uu,\quad U=qtt+2q'tu+q'''uu\quad\text{in}\quad (htt+2itu+kuu)^2
\]

II.\ Entsteht die Form $F$ durch Verdopplung der Form $f$, so
entsteht sie auch durch Verdopplung jeder anderen Form, die mit $f$
in derselben Klasse liegt; mit anderen Worten, die Klasse der Form $F$ entsteht
durch Verdopplung der Klasse der Form $f$ (vgl.\ Art.\ 238). So
entsteht im Beispiel des vorigen Artikels $(5,2,31)$ auch durch Verdopplung
der Form $(11,-5,16)$, die zu $(11,17,40)$ eigentlich äquivalent ist. Aus einer
Klasse, durch deren Verdopplung die Klasse der Form $F$ entsteht,
findet man mithilfe von Problem 260 \textit{alle} solchen Klassen
(falls es mehrere gibt). In unserem Beispiel gibt es keine andere
positive Klasse dieser Art, weil nur eine positive eigentlich primitive
ambige Klasse mit der Determinante $-151$ existiert (nämlich die
Hauptklasse). Aus der Komposition der einzigen negativen ambigen Klasse
$(-1,0,-151)$ mit der Klasse $(11,-5,16)$ entsteht die Klasse $(-11,-5.\ -16)$. Diese ist
also die einzige negative Klasse, durch deren Verdopplung die Klasse
$(5,2,31)$ entsteht.

III.\ Schon die Lösung des Problems im vorigen Artikel zeigt, dass
jede eigentlich primitive (positive) Klasse binärer Formen, die zum
Hauptgeschlecht gehört, durch Verdopplung einer eigentlich primitiven
Klasse derselben Determinante entstehen kann. Der Satz aus Art.\ 261,
der uns die Gewissheit gab, dass \textit{mindestens} der Hälfte aller
für eine gegebene nichtquadratische Determinante
\SourceLeaf{343}
$D$ zuweisbaren Charaktere keine eigentlich primitiven (positiven)
Geschlechter entsprechen können, lässt sich nun so erweitern: \textit{Genau}
der Hälfte aller dieser Charaktere entsprechen tatsächlich solche
Geschlechter; der anderen Hälfte können also keine entsprechen (vgl.\
den Beweis jenes Satzes). In Art.\ 264 wurden alle diese zuweisbaren
Charaktere gleichmäßig in zwei Arten $P$, $Q$ aufgeteilt.
Von den letzteren $Q$ wurde bewiesen, dass sie keinen eigentlich
primitiven (positiven) Formen entsprechen können; bei den übrigen
$P$ blieb dagegen ungewiss, ob jedem tatsächlich stets ein Geschlecht
entspricht. Dieser Zweifel ist nun ganz beseitigt, und wir sind sicher,
dass es in der gesamten Menge der Charaktere $P$ keinen gibt,
dem kein Geschlecht entspricht. --- Daraus folgt auch leicht, dass
bei negativer Determinante in der \textit{negativen} eigentlich primitiven
Ordnung, in der nach Art.\ 264, I alle $P$ unmöglich und \textit{nur}
die $Q$ möglich sind, tatsächlich \textit{alle} $Q$ möglich
sind. Bezeichne nämlich $K$ einen beliebigen Charakter aus $Q$,
$f$ eine beliebige Form aus der negativen eigentlich primitiven
Ordnung der Formen mit der Determinante $D$ und $K'$ ihren
Charakter, so liegt dieser in $Q$. Daraus sieht man leicht, dass
der aus $K$, $K'$ nach Art.\ 246 zusammengesetzte Charakter
zu $P$ gehört. Es gibt also positive eigentlich primitive Formen
mit der Determinante $D$, die ihm entsprechen. Aus der Komposition
einer solchen Form mit $f$ entsteht offenbar eine negative eigentlich
primitive Form mit der Determinante $D$, deren Charakter $K$
ist. --- Auf genau dieselbe Weise zeigt man, dass in der uneigentlich
primitiven Ordnung diejenigen Charaktere, die nach den Regeln aus
Art.\ 264 II, III als \textit{einzig} mögliche gefunden werden,
\textit{alle} möglich sind, ob sie nun $P$ oder $Q$ sind.
--- Diese Sätze gehören, wenn wir uns nicht sehr irren, zu den schönsten
in der Theorie der binären Formen, umso mehr, als sie bei größter
Einfachheit doch so tief verborgen liegen, dass man ohne die Hilfe so
vieler anderer Untersuchungen keinen strengen Beweis für sie geben kann.

\Needspace{8\baselineskip}
\begin{center}
\textit{Theorie der Zerlegung von Zahlen und binären Formen in drei Quadrate.}
\end{center}
Wir gehen nun zu einer anderen Anwendung des vorigen Exkurses über:
der Zerlegung von Zahlen und binären Formen in drei Quadrate. Wir
stellen ihr das folgende Problem voran.

\Needspace{8\baselineskip}
\Article{288}
\textsc{Aufgabe.} \textit{Sei $M$ eine positive Zahl. Gesucht sind
die Bedingungen, unter denen primitive negative binäre Formen mit der
Determinante $-M$ existieren können, die quadratische Reste von
$M$ sind, das heißt, für die $1$ eine charakteristische Zahl ist.}

\SourceLeaf{344}
\textit{Lösung.} Bezeichnen wir mit $\Omega$ die Gesamtheit aller
besonderen Charaktere, die sich aus den Beziehungen der Zahl $1$
sowohl zu den einzelnen (ungeraden) Primteilern von $M$ als auch
zur Zahl $8$ oder $4$ ergeben, wenn diese $M$ teilt.
Diese Charaktere sind offenbar $Rp$, $Rp'$, $Rp''$ usw., wobei
$p$, $p'$, $p''$ usw.\ jene Primteiler bezeichnen; ferner
$1,4$ bei Teilbarkeit durch $4$ und $1,8$, wenn $8$ die Zahl $M$
teilt. Außerdem verwenden wir die Buchstaben $P$, $Q$ im selben
Sinn wie im vorigen Artikel oder in Art.\ 264. Nun unterscheiden wir
die folgenden Fälle.

I.\ Ist $M$ durch $4$ teilbar, so ist $\Omega$ ein vollständiger
Charakter. Nach Art.\ 233 V kann $1$ nur für solche Formen eine
charakteristische Zahl sein, deren Charakter $\Omega$ ist. Nun ist
$\Omega$ offenbar der Charakter der Hauptform $(1,0,M)$, gehört also
zu $P$ und kann daher keiner eigentlich primitiven negativen
Form zukommen. Da es für eine solche Determinante keine uneigentlich
primitiven Formen gibt, gibt es in diesem Fall überhaupt keine primitiven
negativen Formen, die Reste von $M$ sind.

II.\ Ist $M\equiv3\ (\mathrm{mod.}\ 4)$, so gelten genau dieselben Überlegungen, mit der
einzigen Ausnahme, dass in diesem Fall die \textit{uneigentlich}
primitive negative Ordnung existiert. In ihr sind die Charaktere
$P$ möglich oder unmöglich, je nachdem $M\equiv3$ oder $\equiv7\ (\mathrm{mod.}\ 8)$
ist; vgl.\ Art.\ 264, III. Im ersten Fall gibt es also in dieser
Ordnung ein Geschlecht mit dem Charakter $\Omega$, sodass $1$
eine charakteristische Zahl aller darin enthaltenen Formen ist.
Im zweiten Fall kann es überhaupt keine negativen Formen mit dieser
Eigenschaft geben.

III.\ Ist $M\equiv1\ (\mathrm{mod.}\ 4)$, so ist $\Omega$ noch kein vollständiger Charakter;
die Beziehung zur Zahl $4$ muss hinzukommen. Offenbar muss aber
$\Omega$ im Charakter einer Form enthalten sein, deren charakteristische
Zahl $1$ ist; umgekehrt hat jede Form, deren Charakter entweder
$\Omega;\ 1,4$ oder $\Omega;\ 3,4$ ist, die charakteristische Zahl $1$.
Nun ist $\Omega;\ 1,4$ offenbar der Charakter des Hauptgeschlechts, der zu
$P$ gehört und somit in der negativen eigentlich primitiven
Ordnung unmöglich ist. Aus demselben Grund gehört $\Omega;\ 3,4$ zu $Q$
(Art.\ 263), sodass ihm in der negativen eigentlich primitiven Ordnung
ein Geschlecht entspricht, dessen Formen alle die charakteristische
Zahl $1$ haben. Die uneigentlich primitive Ordnung gibt es in
diesem Fall ebenso wenig wie im folgenden.

IV.\ Ist $M\equiv2\ (\mathrm{mod.}\ 4)$, so muss zu $\Omega$ die Beziehung zu $8$
hinzukommen, damit der Charakter vollständig wird, nämlich entweder
$1\mathbin{\textit{und}}3,8$ oder $5\mathbin{\textit{und}}7,8$, wenn $M\equiv2\ (\mathrm{mod.}\ 8)$, und entweder $1\mathbin{\textit{und}}7,8$ oder
$3\mathbin{\textit{und}}5,8$, wenn $M\equiv6\ (\mathrm{mod.}\ 8)$ ist. Im ersten Fall gehört der Charakter
$\Omega;\ 1\mathbin{\textit{und}}3,\allowbreak 8$ offenbar zu $P$, also $\Omega;\ 5\mathbin{\textit{und}}7,\allowbreak 8$ zu $Q$, sodass ihm
\SourceLeaf{345}
ein negatives eigentlich primitives Geschlecht entspricht. Ebenso
gibt es im zweiten Fall ein Geschlecht in der negativen eigentlich
primitiven Ordnung, dessen Formen die vorgeschriebene Eigenschaft
haben, nämlich mit dem Charakter $\Omega;\ 3\mathbin{\textit{und}}5,\allowbreak 8$.

Daraus folgt, dass primitive negative Formen mit der Determinante
$-M$, deren charakteristische Zahl $1$ ist, existieren,
wenn $M$ einer der Zahlen $1$, $2$, $3$, $5$,
$6$ modulo $8$ kongruent ist, und zwar stets in einem einzigen
Geschlecht, das uneigentlich ist, wenn $M\equiv3$. Solche Formen gibt
es überhaupt nicht, wenn $M\equiv0,\allowbreak \ 4\text{ oder }7\ (\mathrm{mod.}\ 8)$. Ist ferner $(-a,-b,-c)$ eine primitive
negative Form mit der charakteristischen Zahl $+1$, so ist
$(a,b,c)$ offenbar eine primitive positive Form mit der charakteristischen
Zahl $-1$. In den ersten fünf Fällen (wenn $M\equiv1,2,3,5,6$) gibt es also
ein primitives positives Geschlecht, dessen Formen die charakteristische
Zahl $-1$ haben, und zwar für $M\equiv3$ ein \textit{uneigentliches}.
In den übrigen drei Fällen (wenn $M\equiv0,4,7$) kann es überhaupt keine
solchen positiven Formen geben.

